\documentclass[10pt,a4paper]{amsart}
\usepackage[margin=19mm]{geometry}
\usepackage{amsmath,amssymb,mathtools}
\usepackage[scr=rsfs,cal=euler]{mathalfa}
\usepackage{xfrac}
\usepackage[unicode,hidelinks,bookmarks=false]{hyperref}
\newcommand{\ie}{i.e.\ }
\newcommand{\cf}{cf.\ }
\newcommand{\resp}{respectively\ }
\newcommand{\Subsection}{\subsection}
\providecommand{\nobreakdash}{\nobreak\hbox{-}}
\setlength{\emergencystretch}{2em}
\multlinegap0pt
\usepackage{bm}
\usepackage{bbm}
\usepackage{dsfont}
\usepackage{xspace}
\usepackage{cancel}
\usepackage{mathtools}
\usepackage{amsthm}
\makeatletter
\@ifundefined{theorem}{\newtheorem{theorem}{Theorem}}{}
\makeatother
\newcommand{\norm}[1]{\left\lVert#1\right\rVert}
\title[Essential norms of pointwise multipliers in the non-algebrai]{Essential norms of pointwise multipliers in the non-algebraic setting}
\author{Tomasz Kiwerski, Jakub Tomaszewski}
\date{Source: arXiv:2212.06723v5 (CC BY 4.0)}
\begin{document}
\maketitle

\noindent\emph{Source and licence.} Essential norms of pointwise multipliers in the non-algebraic setting. Version arXiv:2212.06723v5 (CC BY 4.0), distributed under the Creative Commons Attribution 4.0 International licence, as recorded in the arXiv licence metadata for this exact version and shown on the abstract page https://arxiv.org/abs/2212.06723v5. Full rights notice: licenses/arxiv-2212.06723-CC-BY-4.0.txt.\par\smallskip

\noindent\textit{Editorial scope.} Counted unit: the Theorem whose source label is \texttt{THM: ess norm func + seq} in the article (source0.tex lines 1111--1120, source label \texttt{THM: ess norm func + seq}), delivered together with the article's own complete original proof of it (source0.tex lines 1121--1183, one \texttt{proof} environment of 63 lines). No mathematics is written, repaired or re-derived: statement and proof are the article's own text line for line. Required context retained uncounted: THM: essnorm multipliers function (lines 862--870); THM: Essential norm Kothe sequence (lines 985--992). Every definition, display and label of those lines is retained unchanged. Omitted from this package and not redistributed: 1--861, 871--984, 993--1110, 1184--3483. Nothing inside a retained excerpt is omitted or altered: every retained body is delivered exactly as the article's lines are written, so no omission receipt of any kind accompanies this package. Citations: the retained excerpts carry no citation command at all, so no bibliography entry of the article is transcribed and no reference is redistributed. Reference adaptation: none was needed and none was made; every reference inside the delivered unit resolves inside it, so no unresolved reference or citation is produced. Printed slips kept literal: no printed word, symbol, hypothesis or number of the counted statement or of its proof was altered. Layout and class: the article's own class is \texttt{amsart} with packages including inputenc, babel, fontenc, mathtools, amssymb, amsmath, amsthm, amsfonts, enumerate, csquotes, alphalph, cite, nameref, verbatim; the wrapper is the lane's pinned \texttt{amsart} editorial wrapper at 10pt on A4, which restates the theorem-like environments the retained text needs and adds the article's own macro definitions that the retained text needs (\texttt{\textbackslash norm}), with extra wrapper packages (bm, bbm, dsfont, xspace, cancel, mathtools). The delivered numbering is the unit's own. Assets: no retained span contains a figure environment, a graphics-inclusion command, a PDF-page inclusion, a table or a TikZ picture; no image file is copied into this package, the package declares no asset dependency and no image file is redistributed with it. Licence: version arXiv:2212.06723v5 of this article is distributed under the Creative Commons Attribution 4.0 International licence, as recorded in the arXiv licence metadata for this exact version and shown on the abstract page https://arxiv.org/abs/2212.06723v5. A full rights notice is delivered at licenses/arxiv-2212.06723-CC-BY-4.0.txt. Characters: every retained excerpt is pure ASCII, so the delivered engine, XeLaTeX with its default Latin Modern text font, reports no missing character.
\begin{theorem}[Essential norm of multiplication operators between K{\" o}the function spaces] \label{THM: essnorm multipliers function}
	{\it Let $X$ and $Y$ be two K{\" o}the function spaces with the Fatou property both defined on the same $\sigma$-finite
		measure space $(\Omega,\Sigma,\mu)$. Suppose that either the space $X$ is separable or the space $Y$ is reflexive.
		Then
		\begin{equation*}
			\norm{M_{\lambda} \colon X \rightarrow Y}_{ess} = \norm{\lambda}_{M(X,Y)} = \norm{M_{\lambda} \colon X \rightarrow Y}.
		\end{equation*}
		In particular, there are no non-trivial compact multiplication operators acting from $X$ into $Y$.}
\end{theorem}

\begin{theorem}[Essential norm of multiplication operators between K{\" o}the sequence spaces] \label{THM: Essential norm Kothe sequence}
	{\it Let $X$ and $Y$ be two K{\" o}the sequence spaces with the Fatou property. Suppose that either the space $X$ is reflexive
		or the space $Y$ is separable. Then
	\begin{equation*}
		\norm{M_{\lambda} \colon X \rightarrow Y}_{ess} = \lim\limits_{n \rightarrow \infty} \norm{ \sum_{i=n}^{\infty} \lambda_i e_i }_{M(X,Y)}.
	\end{equation*}
		In particular, the multiplication operator $M_{\lambda} \colon X \rightarrow Y$ is compact if, and only if, $\lambda \in M(X,Y)_o$.}
\end{theorem}

\begin{theorem}[Essential norm of multiplication operators between K{\" o}the spaces] \label{THM: ess norm func + seq}
	{\it Let $X$ and $Y$ be two K{\" o}the spaces with the Fatou property both defined over the same $\sigma$-finite measure space $(\Omega,\Sigma,\mu)$.
		Further, let $\Omega_c$ and $\Omega_a$ denote the non-atomic and purely atomic part of $(\Omega,\Sigma,\mu)$, respectively.
		Suppose that one of the spaces $X$ or $Y$ is reflexive. Then}
		\begin{equation*}
			\norm{M_{\lambda} \colon X \rightarrow Y}_{ess} \approx
				\max \left\{ \norm{M_{\lambda} \colon X(\Omega_c) \rightarrow Y(\Omega_c)}_{ess}, \norm{M_{\lambda} \colon X(\Omega_a) \rightarrow Y(\Omega_a)}_{ess} \right\}
		\end{equation*}
	{\it with an equivalence involving universal constants only.}
\end{theorem}

\begin{proof}
	Without loss of generality, we can assume that the space $M(X,Y)$ is non-trivial.
	Moreover, to make things easier, let us assume that every atom in $\Omega_a$ is equivalent to a singleton.
	Then, since $\Omega = \Omega_c \sqcup \Omega_a$ and the measure space $(\Omega,\Sigma,\mu)$ is assumed to be $\sigma$-finite,
	so $\mu$ restricted to $\Omega_a$ is nothing else but the weighted sum of countably many Dirac measures.
	In other words, $\mu(A) = \mu_c(A) + \sum_{n=1}^{\infty} w_n \delta_{a_n}(A)$ for $A \in \Sigma$,
	where $w_n$'s are positive real numbers and $a_n$'s are atoms in $\Omega_a$.
	
	Take $\lambda \in M(X,Y)$.
	To show the first inequality let $\{K_n\}_{n=1}^{\infty}$ be a sequence of compact operators defined as
	\begin{equation*}
		K_n \colon X \rightarrow Y, \quad
		K_n \colon f \rightsquigarrow \left[ \omega \rightsquigarrow \lambda(\omega) f(\omega) \chi_{\{ a_1,a_2,..., a_{n-1} \}}(\omega) \right],
	\end{equation*}
	where $n \geqslant 2$. We have
	\begin{align*}
		\norm{M_{\lambda} \colon X \rightarrow Y}_{ess} & \leqslant \norm{M_{\lambda} - K_n}_{X \rightarrow Y} \\
			& \leqslant \norm{M_{\lambda} - K_n}_{X(\Omega_c) \rightarrow Y(\Omega_c)} + \norm{M_{\lambda} - K_n}_{X(\Omega_a) \rightarrow Y(\Omega_a)} \\
			& = \norm{M_{\lambda}}_{X(\Omega_c) \rightarrow Y(\Omega_c)} + \norm{\text{Proj}_n M_{\lambda}}_{X(\Omega_a) \rightarrow Y(\Omega_a)}.
	\end{align*}
	Now, it follows from Theorem~\ref{THM: essnorm multipliers function} and Theorem~\ref{THM: Essential norm Kothe sequence} that
	\begin{equation*}
		\norm{M_{\lambda} \colon X(\Omega_c) \rightarrow Y(\Omega_c)}_{ess} = \norm{M_{\lambda}}_{X(\Omega_c) \rightarrow Y(\Omega_c)},
	\end{equation*}
	and, respectively,
	\begin{equation*}
		\norm{M_{\lambda} \colon X(\Omega_a) \rightarrow Y(\Omega_a)}_{ess}
			= \norm{\text{Proj}_n M_{\lambda}}_{X(\Omega_a) \rightarrow Y(\Omega_a)}.
	\end{equation*}
	In consequence, we have
	\begin{equation*}
		\norm{M_{\lambda} \colon X \rightarrow Y}_{ess}
			\leqslant 2\max \left\{ \norm{M_{\lambda} \colon X(\Omega_c) \rightarrow Y(\Omega_c)}_{ess}, \norm{M_{\lambda} \colon X(\Omega_a) \rightarrow Y(\Omega_a)}_{ess} \right\}.
	\end{equation*}

	To prove the second inequality, let us first define two operators
	\begin{equation*}
		\text{Proj}_{\bullet}^X \colon X \rightarrow X(\Omega_{\bullet}), \quad
		\text{Proj}_{\bullet}^X \colon f \rightsquigarrow \left[ \omega \rightsquigarrow f(\omega)\chi_{\Omega_{\bullet}}(\omega) \right]
	\end{equation*}
	and
	\begin{equation*}
		\text{Proj}_{\bullet}^Y \colon Y \rightarrow Y(\Omega_{\bullet}), \quad
		\text{Proj}_{\bullet}^Y \colon f \rightsquigarrow \left[ \omega \rightsquigarrow f(\omega)\chi_{\Omega_{\bullet}}(\omega) \right],
	\end{equation*}
	where $\bullet \in \{a, c\}$.
	It is straightforward to see that $\text{Proj}_{\bullet}^X$ and $\text{Proj}_{\bullet}^Y$ are the norm one projections
	onto $X(\Omega_{\bullet})$ and, respectively, $Y(\Omega_{\bullet})$.
	Now, we have
	\begin{align*}
		\norm{M_{\lambda} \colon X(\Omega_{\bullet}) \rightarrow Y(\Omega_{\bullet})}_{ess}
			& = \inf \left\{ \norm{M_{\lambda} - K}_{X(\Omega_{\bullet}) \rightarrow Y(\Omega_{\bullet})} \colon K \in \mathscr{K}(X(\Omega_{\bullet}),Y(\Omega_{\bullet})) \right\} \\
			& = \inf \left\{ \norm{\text{Proj}_{\bullet}^Y(M_{\lambda} - K)\text{Proj}_{\bullet}^X}_{X \rightarrow Y} \colon K \in \mathscr{K}(X,Y) \right\} \\
			& \leqslant \inf \left\{ \norm{M_{\lambda} - K}_{X \rightarrow Y} \colon K \in \mathscr{K}(X,Y) \right\} \\
			& = \norm{M_{\lambda} \colon X \rightarrow Y}_{ess}.
	\end{align*}
	Thus,
	\begin{equation*}
		\max \left\{ \norm{M_{\lambda} \colon X(\Omega_c) \rightarrow Y(\Omega_c)}_{ess}, \norm{M_{\lambda} \colon X(\Omega_a) \rightarrow Y(\Omega_a)}_{ess} \right\}
			\leqslant \norm{M_{\lambda} \colon X \rightarrow Y}_{ess}.
	\end{equation*}
	This ends the proof.
\end{proof}

\end{document}
